\section{Data and names}
\label{sec:qs-data}

All conventions are those of the manuscript, Chapters 1--11 (Sections 11.1--11.5 final), of the two-letter theory (\path{step2/twoletter/theory/twoletter_theory.tex}) and of the composite theory (\path{step2/composite_qsk/composite/theory/composite_theory.tex}). References of the form Def c8-K, Prop c11-oddbound are to labels of the manuscript. Weight $e^{+S/\hbar}$, $S=\tfrac12\langle e,Qe\rangle+I$; on arity zero $s=Q+\{I,-\}$; edges $P=P(\varepsilon,T)=\int_\varepsilon^T(Q^{GF}\otimes1)K_t\,dt$ (Def c8-PT) and $P\star=\int_\varepsilon^TQ^{GF}e^{-tH}dt$ the operator it represents; the special edge is $K_\varepsilon-K_0=-\int_0^\varepsilon(H\otimes1)K_t\,dt$, the factor $e^{-\varepsilon|p|^2}-1$ times the components of $K_0$ (Prop c4-defect1). No dimensional regularisation enters in any role; the only infrared datum is the transfer endpoint $T$, which is set to $\infty$ wherever the graph sums converge at generic momenta (Prop c8-Tsafe) and kept finite along $K_b$ (Def c8-K). The unit is
\[
u=\frac{\hbar\operatorname{tr}_{\mathfrak g}(\mathrm{ad}_X\mathrm{ad}_Y)}{16\pi^2(X,Y)}=-\frac{g^2N}{32\pi^2}\qquad(\mathfrak g=\mathfrak{su}(N)).
\]

\textbf{Generators and components.} As in Chapter 8, $\xi\in\Sigma$ with two-component data $(\xi_\alpha{}^i,\bar\xi_{\dot\alpha i})$ (generators $Q^i_\alpha$, $\bar Q_{\dot\alpha i}$), $\eta\in\Sigma^*$ with data $(\eta^\alpha{}_i,\bar\eta^{\dot\alpha i})$ (generators $S^\alpha{}_i$, $\bar S^{\dot\alpha i}$), $K_b$ the special conformal field $v(x)=(x,x)b-2(b,x)x$, $P_b$ the translation, $\Delta$ the dilatation ($v=2E$), $M_\omega$ the rotation of $V$ with the antisymmetric endomorphism $\omega$. The letters $Q_1,S_1,D_1,K_1,P_1,M_1,R_1$ are the one-loop components $\hbar L_1(O)^{(1)}(x)$ of Def c8-components; $Q$ alone is the linear BV differential. The $\mathrm{ad}_\Delta$ eigenvalue $\lambda_x$ is $-1$ on $\Sigma$, $+1$ on $\Sigma^*$, $+2$ for $K_b$, $-2$ for $P_b$, $0$ for $\Delta$ and the rotations; the degree of the Killing data of $x$ in the position is $0$, $1$, $2$, $0$, $1$, $1$ (Section 10.9).

\textbf{Letters and pairs.} Letters, Weyl weights, degrees, the pair weight $w_{X_a}+w_{X_b}$, the pair degree and the Koszul signs are those of Def c10-letters and \eqref{eq:qs-decomp} below, which is eq c10-decomp. The slots are $c$, $\AV$, $\Ap:=A|_{W_{\mathbb R}^\perp}$, $\lambda_+:=\lambda|_{S_+\otimes W}$, $\lambda_-:=\lambda|_{S_-\otimes W^*}$, $A^*|_V$, $A^*|_\perp$, $\lambda^*_+:=\lambda^*|_{S_+\otimes W^*}$, $\lambda^*_-:=\lambda^*|_{S_-\otimes W}$, $c^*$, with Weyl weights $0,1,1,\tfrac32,\tfrac32,3,3,\tfrac52,\tfrac52,4$ and degrees $1,0,0,0,0,-1,-1,-1,-1,-2$; two-component $c$, $A_m$, $A_{ij}$, $\lambda_\alpha{}^i$, $\bar\lambda_{\dot\alpha i}$, $A^*_m$, $A^*_{ij}$, $\lambda^{*\alpha}{}_i$, $\bar\lambda^{*\dot\alpha i}$, $c^*$ (Def c8-twocomp). The arity-zero counterterm of a local functional $M$ is written
\[
\partial_\delta I^{CT}[\varepsilon](M)=\sum_{m\ge1}\varepsilon^{-m}\,\mathrm{ct}_m(M)+\log\varepsilon\;\mathrm{lc}(M),
\]
with $\mathrm{ct}_m$, $\mathrm{lc}$ two-body by Thm c10-decomp; the arity-one counterterm of $M$ along $x$ (the singular part of the cycles through two letters of $M$ with one vertex $S^{(1)}(x)$ on the cycle, Def c11-graphs for $x\in\{\Delta,P_b\}$, Definition \ref{def:qs-ct1} below in general) is written $\mathrm{ct}^{(1)}_x(M)=\sum_m\varepsilon^{-m}\mathrm{ct}^{(1)}_{x,m}(M)+\log\varepsilon\,\mathrm{lc}^{(1)}_x(M)$.

\begin{definition}[The terms of $S^{(1)}(x)$]
\label{def:qs-terms}
For $\xi\in S_+\otimes W$ the functional $S^{(1)}(\xi)=\int\mathrm{dvol}\big(-(\Gamma(\xi,\lambda),A^*)+\tfrac12(\rho(F_A)\xi,\lambda^*)\big)$ (eq c2-S1Q) is the sum of the seven terms of Table \ref{tab:qs-terms}: four quadratic terms $T_1,\dots,T_4$ and three cubic terms $T_5$. Each term is a vertex with the listed slots; $T_3$, $T_4$ carry one derivative on the leg $A$. For $\xi\in S_-\otimes W^*$ the table holds with $\lambda_\pm$, $\lambda^*_\pm$ exchanged. For $\eta\in\Sigma^*$,
\begin{equation}
S^{(1)}(\eta)=\sum_\mu\int x^\mu\,s_\mu+\int s_0,\qquad \int s_\mu:=S^{(1)}(\xi_\mu),\quad \xi_\mu:=\rho(e_\mu)\eta\in\Sigma,\qquad s_0:=-\big(\rho(A|_\perp)\eta,\lambda^*\big),
\label{eq:qs-Seta}
\end{equation}
(eq c2-S1S, Lemma c11-lines (ii)), with $e_\mu$ an orthonormal basis of $V$; for $\eta\in S_+\otimes W^*$ the spinors $\xi_\mu$ lie in $S_-\otimes W^*$ and $s_0$ has the slots $(\Ap,\lambda^*_+)$, for $\eta\in S_-\otimes W$ they lie in $S_+\otimes W$ and $s_0$ has the slots $(\Ap,\lambda^*_-)$. For even $x$ with vector field $v$, $S^{(1)}(x)=\int\big((v.A,A^*)-(v.\lambda,\lambda^*)-(v.c,c^*)\big)$ (Def c7-vdot), quadratic, with the five slot pairs $(\AV,A^*|_V)$, $(\Ap,A^*|_\perp)$, $(\lambda_\pm,\lambda^*_\pm)$, $(c,c^*)$.
\end{definition}

\begin{longtable}{l|p{3.6cm}|p{7.2cm}|p{3.2cm}}
term & index-free & two-component, $\xi_\alpha{}^i=u_\alpha z^i$ (Rem c8-twist) & slots\\
\hline
$T_1$ & $-(\Gamma(\xi,\lambda)|_V,A^*|_V)$ & $-i\big(\bar\lambda_{\dot\alpha i}z^i\bar\sigma^{m\dot\alpha\alpha}u_\alpha,A^*_m\big)$ & $\lambda_-$, $A^*|_V$\\
$T_2$ & $-(\Gamma(\xi,\lambda)|_\perp,A^*|_\perp)$ & $\tfrac12\varepsilon_{jikl}z^j\big((u\lambda^i),A^{*kl}\big)$ & $\lambda_+$, $A^*|_\perp$\\
$T_3$ & $\tfrac12(\rho(dA|_V)\xi,\lambda^*)$ & $-\tfrac14\big(F^0_{mn}(\sigma^m\bar\sigma^nu)_\alpha,\lambda^{*\alpha}{}_iz^i\big)$, $F^0=dA|_V$ & $\AV$ (one derivative), $\lambda^*_+$\\
$T_4$ & $\tfrac12(\rho(dA|_\perp)\xi,\lambda^*)$ & $-\tfrac i2\big(u^\alpha\sigma^m_{\alpha\dot\alpha}\partial_m(A_{ij}z^j),\bar\lambda^{*\dot\alpha i}\big)$ & $\Ap$ (one derivative), $\lambda^*_-$\\
$T_{5}$ & $\tfrac12(\rho(\tfrac12[A\wedge A])\xi,\lambda^*)$ & the commutator terms of $F_{mn}$, $D_m(A_{ij}z^j)$ and $-\tfrac14\big([A^{ik},A_{kj}z^j]u_\alpha,\lambda^{*\alpha}{}_i\big)$ & $(\AV,\AV,\lambda^*_+)$, $(\AV,\Ap,\lambda^*_-)$, $(\Ap,\Ap,\lambda^*_+)$\\
$s_0$ & $-(\rho(A|_\perp)\eta,\lambda^*)$ & $\big(A^{kj}\varepsilon_{\alpha\beta}\eta^\beta{}_j,\lambda^{*\alpha}{}_k\big)-\big(A_{kj}\varepsilon_{\dot\beta\dot\alpha}\bar\eta^{\dot\alpha j},\bar\lambda^{*\dot\beta k}\big)$ & $(\Ap,\lambda^*_\pm)$\\
\end{longtable}
\begin{center}Table \refstepcounter{table}\thetable\label{tab:qs-terms}: the terms of $S^{(1)}(\xi)$, $\xi\in S_+\otimes W$, and the intrinsic term of $S^{(1)}(\eta)$. The two-component forms are those of Rem c8-twist for $\xi_\alpha{}^i=u_\alpha z^i$, $\bar\xi=0$, and of Def c8-twocomp for $\rho(A|_\perp)\eta$; for general $\xi$ replace $u_\alpha z^i$ by $\xi_\alpha{}^i$.\end{center}

The census of Section \ref{sec:qs-census} uses only the slot structure of Table \ref{tab:qs-terms} and of Def c1-vertex; the engines use the explicit realisation of $S^{(1)}$ in \path{superconformal/model.py} (\texttt{S1\_odd}, \texttt{S1\_even}), which is the realisation of Chapter 2.

\begin{definition}[Defect cycles and two-body kernels]
\label{def:qs-defect}
Let $x$ be a generator, $X_a$, $X_b$ two letters of any slots, $0<\varepsilon<T\le\infty$.
\begin{itemize}
\item (i) A defect cycle of the pair $(X_a,X_b)$ along $x$ is a chain $X_a-e_0-V_1-e_1-\cdots-V_n-e_n-X_b$ (Def c10-chain) whose vertices $V_j$ are vertices of $I$ and exactly one vertex $V_s$ given by a term of $S^{(1)}(x)$, entered and left through two of its slots (the third slot of a cubic term is an output), and whose edges are components of $P(\varepsilon,T)$ and exactly one special edge $e_\sigma=K_\varepsilon-K_0$, with the components of $K_0$ (a field slot with its antifield slot, either one possibly a slot of a letter). The vertex $V_s$ is placed in one of three ways: \emph{on a line}, with both legs on edges $P$; \emph{at the special edge}, with one leg on $e_\sigma$; \emph{cubic}, a cubic term with its two cycle legs on edges $P$ or one of them on $e_\sigma$.
\item (ii) $\mathcal D_x(X_a,X_b;\varepsilon,T)$ is the sum of the one-loop graphs of the flow of the defect vertex of Lemma c8-transfer, applied to $e^{(J+\delta M(\varepsilon))/\hbar}$ with $M=\int\omega\,X_aX_b$, that are $\delta$-linear, of arity one along $x$, and whose cycle is a defect cycle of the pair, amputated at the outputs, the outputs left open as in Def c10-pair. The signs are those of the flow: the edge rule of Prop c3-graphs, the signs of Lemma c10-sign along the chain from $X_a$, and the sign of the special edge of Prop c11-master ($\{f,g\}_\varepsilon-\{f,g\}$ is minus $f$ and $g$ joined by $K_\varepsilon-K_0$). Operationally, a chain read from $X_a$ with $V_s$ at position $s$ and the special edge $e_\sigma$ carries, beyond the sign of Lemma c10-sign, the factor $-(-1)^{n(<\sigma)+|S^{(1)}(x)|\,(n(<s)+[s<\sigma])}$, with $n(<j)$ the number of odd outputs before position $j$, $[s<\sigma]=1$ if $V_s$ precedes $e_\sigma$, and $|S^{(1)}(x)|$ the total parity of the functional $S^{(1)}(x)$: $0$ for odd $x$, $1$ for even $x$, for which this is the factor of Def c11-sp up to the overall sign of the special edge; the engines fix the rule in code against the gates G-pole and G-reg of Section \ref{sec:qs-engines}.
\item (iii) The finite special vertex of the pair is
\begin{equation}
\mathcal F_x(X_a,X_b):=-\sum_{m\ge1}\frac{(-1)^m}{m!}\Big(\big\{I,\mathrm{ct}^{(1)}_{x,m}(X_aX_b)\big\}_{H^m}+\big\{S^{(1)}(x),\mathrm{ct}_m(X_aX_b)\big\}_{H^m}\Big),
\label{eq:qs-finitesp}
\end{equation}
with $\{-,-\}_{H^m}$ the bracket built from $(H^m\otimes1)K_0$ in place of $K_0$ (eq c10-finitesp).
\item (iv) The two-body kernel of $(X_a,X_b)$ along $x$ is
\begin{equation}
\Upsilon_x(X_a,X_b):=\big[\varepsilon^0\big]\,\mathcal D_x(X_a,X_b;\varepsilon,T)+\mathcal F_x(X_a,X_b),
\label{eq:qs-kernel}
\end{equation}
the coefficient of $\varepsilon^0$ in the expansion of Theorem \ref{thm:qs-simplex}, a local density of the outputs at the point of the pair with values in the tensor product of the fibres and colour spaces of the two letters.
\end{itemize}
For $x\in\{\Delta,P_b\}$ and the placement on a line, $\mathcal D_x$ restricted to these cycles is the special-edge composite cycle $Z_x$ of Def c11-sp.
\end{definition}
